\section{Methodology}
\label{sec:method}

\subsection{Profile Coherence}
\label{sec:method:pck}

\subsubsection{Ordinal band scale}
Customers and assets share the same 4-point ordinal scale, running from Conservative (0) through Income (1) and Balanced (2) to Aggressive (3). The customer band $b_u$ is taken from FAR-Trans's \texttt{riskLevel} field. Both declared values and regression-imputed \texttt{Predicted\_*} entries are mapped to the same ordinal scale; customers whose entry is \texttt{Not\_Available} or missing are excluded from all coherence computations. The asset band $b_i$ follows the hierarchical rule in Section~\ref{sec:setup:bandassignment}.

\subsubsection{Pairwise discordance}
\begin{equation}
d(u, i) = |b_u - b_i| \in \{0, 1, 2, 3\}.
\end{equation}
A recommendation is \emph{profile-coherent} iff $d \leq 1$. The tolerance of $\pm 1$ band reflects MiFID II's adjacency allowance for borderline suitability; Section~\ref{sec:limitations} discusses this design choice.

\subsubsection{Aggregation}
\begin{equation}
\text{PC@}k(u) = \frac{1}{k} \left|\{ i \in \text{top}_k(u) : d(u, i) \leq 1 \}\right|.
\end{equation}
Two implementation rules complete the definition. First, recommenders returning fewer than $k$ items use the actual number of items returned as the denominator. Second, customers with $b_u = \emptyset$ are dropped from the aggregate. Per-split PC@$k$ is the unweighted mean across eligible customers and the headline numerics average over all splits.

\subsubsection{Band-conditional random baseline}
The closed-form expectation of PC@$k$ under a uniformly random recommender for band $b$ is given by
\begin{equation}
\pi(b) = \frac{|\{ i \in A : |b - b_i| \leq 1 \}|}{|A|},
\end{equation}
where $A$ is the asset universe. Section~\ref{sec:setup:bandassignment} reports the FAR-Trans asset-band distribution and the resulting $\pi(b)$ values.

\subsubsection{PC-lift@$k$}
Per-user PC@$k$ is normalised against the random reference at the customer's own band:
\begin{equation}
\text{PC-lift@}k(u) = \frac{\text{PC@}k(u)}{\pi(b_u)}.
\end{equation}
Lift $> 1$ beats the band-conditional random rate, while lift $= 1$ matches it. Because 77.9\% of FAR-Trans assets sit in the Conservative, Income, and Balanced bands, a recommender that pushes Income assets at everyone would score well on raw PC@$k$ for any user whose tolerance set includes Income. Dividing by $\pi(b_u)$ strips the band-specific easiness of the asset universe and makes per-band values commensurate.

\subsection{Stratified, Profile-Coherent LightGCN}
\label{sec:method:pclgcn}

\subsubsection{Backbone}
LightGCN is the backbone: a differentiable loss with per-customer embeddings lets us attach the coherence margin to the BPR objective. Random Forest is not extended as its objective uses asset-level technical indicators only and carries no customer signal.

\subsubsection{Stratified architecture}
We train four band-stratified sub-models, one per declared MiFID risk band. Each sub-model's training graph contains only the interactions of customers assigned to that band. At inference, each customer is routed to the sub-model matching their declared band; customers with no declared band default to the Balanced sub-model.

\subsubsection{Profile-coherent margin loss}
BPR samples a positive item from observed Buys and a negative item uniformly. Independently of the BPR sampling, we add a second margin term that samples one coherent asset $i_c$ (satisfying $d(u, i_c) \leq 1$) and one discordant asset $i_d$ (satisfying $d(u, i_d) > 1$) from the sub-model's asset universe, and pushes the coherent score above the discordant score:
\begin{equation}
\mathcal{L}_{\text{coh}}(u) = \text{softplus}\!\left(s(u, i_d) - s(u, i_c)\right),
\end{equation}
where $s(u, i)$ is the LightGCN inner-product score between customer $u$ and asset $i$.

\subsubsection{Total loss}
\begin{equation}
\mathcal{L} = \mathcal{L}_{\text{BPR}} + \omega \|\Theta\|_2^2 + \lambda \cdot \mathcal{L}_{\text{coh}},
\end{equation}
with $\omega$ the L2 weight decay on embedding parameters $\Theta$ and mixing weight $\lambda \geq 0$. BPR remains the primary ranking signal and only the treatment trial activates the coherence term.

\subsubsection{Ablation Setup}
$\lambda = 0$ leaves stratified BPR alone, isolating the architecture's contribution, while $\lambda = 1$ adds the coherence margin. The two trials decompose the full intervention into stratification and PC-loss pieces.
